% titel: Blatt 10, A6: Schallpegel
% skript: Kap. Akustik, Schalldruckpegel und Addition inkohärenter Quellen

% Lautstaerke als dB, 10 log10(I/I0); N gleichlaut inkohaerente Quellen
% addieren Intensitaet linear -> +10 log10(N) dB.

Schall wird nicht in linearer Intensität, sondern logarithmisch als
\emph{Schallpegel} angegeben. Der Grund liegt in der riesigen Spanne h"orbaren
Schalls: Die leiseste h"orbare Intensität ($I_0\approx\qty{1e-12}{W/m^2}$)
und die Schmerzgrenze unterscheiden sich um zw"olf Zehnerpotenzen. Ein
logarithmisches Ma"ss bringt diese Spanne auf eine handliche Skala.

\begin{werkzeugkasten}[Schallpegel: Dezibel-Skala]
\begin{equation}
  L=10\log_{10}\frac{I}{I_0}\ \mathrm{dB},\qquad
  I_0=\qty{1e-12}{W/m^2}\ \text{(H"orschwelle)}.
\end{equation}
Der Faktor 10 steckt drin, weil 1~dB ``gerade noch h"orbarer Unterschied''
ist: 10~dB bedeuten eine Verzehnfachung der Intensit"at, 20~dB eine
Verhundertfachung usw. Ein \emph{Verh"altnis} der Intensit"aten wird also nur
vom Logarithmus des Verh"altnisses bestimmt; f"ur zwei Pegel $L_1,L_2$
\begin{equation}
  L_2-L_1=10\log_{10}\frac{I_2}{I_1}\ \mathrm{dB}.
\end{equation}
Wichtig ist die \emph{Logarithmus-Logik}: eine Verdopplung der Intensit"at
ist nur $+3$~dB ($10\log_{10}2\approx3$) --- nicht h"obstens laut, sondern
``etwas h"orbar lauter''.
\end{werkzeugkasten}

\begin{werkzeugkasten}[Warum N gleichlaut Quellen nur $+10\log_{10}N$ dB]
Zwei unabh"angige, \emph{inkoh"arente} Schallquellen (z. B. zwei Trillerpfeifen,
zwei Trommeln) haben keine feste Phasenbeziehung. Mittelt sich die
Interferenz heraus, so \emph{addieren sich die Intensit"aten}: $I_N=N\,I_1$.
Damit
\begin{equation}
  L_N=10\log_{10}\frac{N I_1}{I_0}
  =10\log_{10}\frac{I_1}{I_0}+10\log_{10}N
  =L_1+10\log_{10}N.
\end{equation}
Zahlen:
\begin{itemize}
  \item 10 gleiche Quellen: $+10\log_{10}10=+10$~dB.
  \item 100 Quellen: $+20$~dB; 1000 Quellen: $+30$~dB.
  \item 2 Quellen: nur $+3$~dB --- die Logarithmus-Skala macht Schall
        ``langsam laut''.
\end{itemize}
Das ist der physikalische Kern beider Teilaufgaben: die Antwort h"angt nur
vom \emph{Verh"altnis} der Anzahlen ab, also direkt vom $10\log_{10}N$.
\end{werkzeugkasten}

\begin{werkzeugkasten}[Schmerzgrenze und Grössenordnungen]
Richtwerte in geringer Entfernung: Flüstern ~30~dB, normale Unterhaltung
~60~dB, laute Stra"se ~80~dB, Rockkonzert/Diskothek ~100~dB,
Trillerpfeife in geringer Entfernung ~120~dB, \textbf{Schmerzschwelle
$\approx130$~dB}.
Ein Zuwachs von 10~dB wird als ``grob doppelt so laut'' empfunden --- wegen
der Logarithmus-Skala psychologisch etwa linear zur dB-Zahl.
\end{werkzeugkasten}

\begin{achtung}[Pegel vs.\ Intensit"at bei mehreren Quellen]
\textbf{Bei inkoh"arenten Quellen Intensit"aten addieren, nicht Drücke} und
\textbf{nicht direkt dB}: $I_{\text{tot}}=\sum I_i$. Nur bei \emph{exakt
gleichen, koh"arenten} Quellen (gleiche Phase) w"urden die \emph{Amplituden/
Drücke} addieren, was doppelt so viele Amplituden $\Rightarrow 4\times$ die
Intensit"at $\Rightarrow +6$~dB w"are. Bei verschiedenen Quellen (Pfeifen,
Trommeln) gilt die inkoh"arente Regel $+10\log_{10}N$.
\end{achtung}

\begin{achtung}[F"ulle beim Umkehren]
Sucht man die Anzahl $N$ zu einer gewünschten Pegel-Differenz
$\Delta L=L_2-L_1$, löst man
\begin{equation}
  N=10^{\Delta L/10}\;.
\end{equation}
F"ur $\Delta L=30$~dB ist $N=10^{3}=1000$; f"ur $\Delta L=10$~dB
$N=10$. Nicht vergessen, dass ``\emph{mehr als}`` bedeuten kann, dass die
errechnete Zahl allein noch nicht reicht (Achtung auf $\ge$ vs.\ $>$).
\end{achtung}
